\chapter{Реализация модуля сбора и потоковой обработки сетевых событий}

\section{Модуль сбора данных}

Модуль сбора данных реализован в файле \texttt{src/collector/packet\_capture.py} и содержит два основных класса: \texttt{PacketCapture} и \texttt{PacketGenerator}.

\subsection{Класс PacketCapture}

Класс \texttt{PacketCapture} обеспечивает захват сетевого трафика из двух источников:

\begin{enumerate}
    \item \textbf{Live-захват} — перехват пакетов через сетевой интерфейс с использованием библиотеки Scapy. Захват выполняется в отдельном потоке, чтобы не блокировать основной поток программы.
    \item \textbf{Загрузка PCAP-файлов} — чтение ранее сохранённых дампов трафика через функцию \texttt{rdpcap()}.
\end{enumerate}

Для каждого захваченного пакета выполняется извлечение базовых признаков (см. листинг~\ref{lst:extract_features}):

\begin{itemize}
    \item \textbf{IP-адреса} источника и назначения (src\_ip, dst\_ip)
    \item \textbf{Протокол} транспортного уровня (TCP=6, UDP=17, ICMP=1)
    \item \textbf{Размер пакета} в байтах (length)
    \item \textbf{TTL} (Time To Live) — время жизни пакета
    \item \textbf{Порты} источника и назначения (src\_port, dst\_port) — для TCP/UDP
    \item \textbf{TCP-флаги} (SYN, ACK, FIN, RST) — для TCP-пакетов
    \item \textbf{Размер TCP-окна} (tcp\_window)
\end{itemize}

\begin{lstlisting}[caption=Извлечение признаков из пакета, label=lst:extract_features, language=Python]
def _extract_features(self, packet):
    if not packet.haslayer(IP):
        return None
    ip_layer = packet[IP]
    features = {
        'timestamp': packet.time if hasattr(packet, 'time') else time.time(),
        'src_ip': ip_layer.src,
        'dst_ip': ip_layer.dst,
        'proto': ip_layer.proto,
        'length': len(packet),
        'ttl': ip_layer.ttl,
    }
    if packet.haslayer(TCP):
        tcp = packet[TCP]
        features.update({
            'src_port': tcp.sport,
            'dst_port': tcp.dport,
            'flags': str(tcp.flags),
            'tcp_window': tcp.window,
        })
    elif packet.haslayer(UDP):
        udp = packet[UDP]
        features.update({
            'src_port': udp.sport,
            'dst_port': udp.dport,
        })
    elif packet.haslayer(ICMP):
        icmp = packet[ICMP]
        features.update({
            'icmp_type': icmp.type,
            'icmp_code': icmp.code,
        })
    return features
\end{lstlisting}

\subsection{Класс PacketGenerator}

Класс \texttt{PacketGenerator} предназначен для генерации синтетических сетевых данных, что позволяет тестировать платформу без необходимости доступа к реальному сетевому трафику. Реализована поддержка четырёх типов аномального трафика:

\subsubsection{DDoS-атака (Distributed Denial of Service)}

Генерируется большое количество SYN-пакетов с множества различных IP-адресов на 80-й порт целевого хоста. Характеризуется высокой частотой пакетов и малым размером:

\begin{itemize}
    \item Пакеты отправляются с интервалом 0.5 мс
    \item Размер пакета: 40--80 байт
    \item Флаг TCP: SYN
    \item IP-адреса источника: диапазон 10.0.0.0/24
\end{itemize}

\subsubsection{Сканирование портов (Port Scan)}

Генерируются одиночные SYN-пакеты с одного источника на различные порты назначения целевого хоста. Характеризуется большим количеством уникальных портов назначения:

\begin{itemize}
    \item Пакеты отправляются с интервалом 0.3 мс
    \item Размер пакета: 60 байт (фиксированный)
    \item Порты назначения: случайные в диапазоне 1--65535
\end{itemize}

\subsubsection{Подбор паролей (Brute Force)}

Генерируются TCP-соединения на 22-й порт (SSH) с одного IP-адреса. Характеризуется регулярными попытками подключения с различными комбинациями флагов:

\begin{itemize}
    \item Интервал между попытками: 20 мс
    \item Размер пакета: 100--300 байт
    \item Различные комбинации TCP-флагов (PA, A, FA, R)
\end{itemize}

\subsubsection{Утечка данных (Data Exfiltration)}

Генерируются крупные TCP-пакеты на внешний IP-адрес через 443-й порт (HTTPS). Характеризуется большим размером пакетов:

\begin{itemize}
    \item Размер пакета: 800--1500 байт (среднее 1200)
    \item Высокая частота отправки
    \item IP назначения: внешний адрес (203.0.113.0/24)
\end{itemize}

\section{Модуль потоковой обработки}

Модуль обработки реализован в файле \texttt{src/processing/feature\_extractor.py} и содержит класс \texttt{FeatureExtractor}.

\subsection{Агрегация во временные окна}

Основной метод \texttt{aggregate\_window()} выполняет группировку сырых пакетов во временные окна фиксированного размера (по умолчанию 2 секунды). Для каждого окна вычисляется 23 признака (см. таблицу~\ref{tab:features}).

\begin{table}[H]
\centering
\caption{Признаки, вычисляемые для каждого временного окна}
\label{tab:features}
\begin{tabular}{p{4cm}p{4.5cm}p{6cm}}
\toprule
\textbf{Группа} & \textbf{Признаки} & \textbf{Описание} \\
\midrule
Количественные & packet\_count, bytes\_total, bytes\_mean, bytes\_std, bytes\_min, bytes\_max & Статистики по количеству и размеру пакетов \\
\addlinespace
Сетевые & unique\_src\_ip, unique\_dst\_ip, unique\_src\_port, unique\_dst\_port & Количество уникальных адресов и портов \\
\addlinespace
Протокольные & proto\_tcp, proto\_udp, proto\_icmp, syn\_count, syn\_ratio & Распределение по протоколам, количество SYN \\
\addlinespace
Временные & packets\_per\_sec, bytes\_per\_sec & Интенсивность трафика \\
\addlinespace
Энтропийные & src\_entropy, dst\_port\_entropy & Разнообразие IP и портов \\
\bottomrule
\end{tabular}
\end{table}

\subsection{Энтропийные метрики}

Особое внимание уделено энтропийным метрикам, которые позволяют количественно оценить разнообразие IP-адресов и портов во временном окне. Энтропия вычисляется по формуле:

\begin{equation}
\label{eq:entropy}
H = -\sum_{i=1}^{n} p_i \log_2(p_i + \varepsilon)
\end{equation}

где:
\begin{itemize}
    \item \(n\) — количество уникальных значений признака в окне;
    \item \(p_i\) — доля i-го значения;
    \item \(\varepsilon = 10^{-10}\) — малая константа для избежания логарифма нуля.
\end{itemize}

Высокая энтропия IP-адресов источника (src\_entropy) может указывать на DDoS-атаку, а высокая энтропия портов назначения (dst\_port\_entropy) — на сканирование портов.

\begin{lstlisting}[caption=Вычисление энтропии, label=lst:entropy, language=Python]
def _entropy(self, series):
    counts = series.value_counts()
    probs = counts / counts.sum()
    return -np.sum(probs * np.log2(probs + 1e-10))
\end{lstlisting}

\subsection{Реализация оконной агрегации}

Процесс агрегации включает следующие шаги:

\begin{enumerate}
    \item Преобразование временных меток в числовой формат и сортировка.
    \item Вычисление номера окна для каждого пакета:
    \[
    \text{window\_idx} = \left\lfloor \frac{t - t_{\min}}{\Delta t} \right\rfloor
    \]
    где \(t\) — временная метка пакета, \(t_{\min}\) — минимальное время в наборе, \(\Delta t\) — размер окна.
    \item Группировка пакетов по номеру окна.
    \item Для каждой группы вычисление всех 23 признаков.
\end{enumerate}

\begin{lstlisting}[caption=Метод агрегации окон, label=lst:aggregation, language=Python]
def aggregate_window(self, df):
    if df.empty:
        return pd.DataFrame()
    df = df.copy()
    df['timestamp'] = pd.to_numeric(df['timestamp'], errors='coerce')
    df = df.dropna(subset=['timestamp'])
    df = df.sort_values('timestamp')
    start_time = df['timestamp'].min()
    df['window'] = ((df['timestamp'] - start_time)
                    // self.window_size).astype(int)
    agg_features = []
    for win, group in df.groupby('window'):
        features = {
            'window': win,
            'window_start': start_time + win * self.window_size,
            'packet_count': len(group),
            'bytes_mean': group['length'].mean(),
            'src_entropy': self._entropy(group['src_ip']),
            # ... и другие признаки
        }
        agg_features.append(features)
    return pd.DataFrame(agg_features)
\end{lstlisting}
